To evaluate the behavior of each transform within our compression pipeline, we
report several standard metrics drawn from image coding. These metrics quantify
both reconstruction fidelity and the effectiveness of our bitstream reduction
steps.

\paragraph{Mean Squared Error (MSE).}
Given an original image $I$ and its reconstruction $\hat{I}$, the MSE is
\begin{equation}
    \mathrm{MSE}
    = \frac{1}{N} \sum_{i=1}^{N} (I_i - \hat{I}_i)^2,
\end{equation}
where $N$ is the number of pixels. Lower MSE indicates a reconstruction closer
to the original.

\paragraph{Peak Signal-to-Noise Ratio (PSNR).}
PSNR provides a logarithmic measure of fidelity, defined as
\begin{equation}
    \mathrm{PSNR}
    = 10 \log_{10}\!\left(\frac{255^2}{\mathrm{MSE}}\right) \ \mathrm{dB},
\end{equation}
assuming 8-bit images. Higher PSNR corresponds to higher reconstruction
quality. We use PSNR as the primary distortion metric in all rate–distortion
plots.

\paragraph{Compression Ratio (CR).}
As a reminder, we define the compression ratio as
\[
\mathrm{CR} = \frac{\text{uncompressed size}}{\text{compressed size}}.
\]
Unless otherwise noted, all CR values in the rate--distortion plots refer to the 
final bitstream after entropy coding. For completeness, we also report the
pre-entropy (“direct”) CR in one diagnostic figure to examine how much gain each
transform receives from the entropy coder.

We will also continue to use Shannon's definition of entropy discussed in Section \ref{sec:lit review}. It gives an accurate picture of how well each transform successfully decorrelated the input image.

\medskip
Together, these metrics provide the basis for the rate–distortion evaluations
reported in Section \ref{sec:results}.